\section{Analyticity along extreme-point motions}
\label{sec:physical-ward-analytic-continuation}

To continue a Ward identity away from a line, we need analyticity along
certain one-coordinate motions. Smoothness of the hierarchy does not
suffice. We derive the required analyticity from reflection positivity,
positive-time heat evolution, and the collinear angular identities. The
argument uses only finite-dimensional exponentials of individual spin
operators; an action of the full rotation group is not an input.

Let $S$ be the smooth, translation-invariant, reflection-positive hierarchy
already constructed, with finite-dimensional filtered field spaces $E_d$.
Write $D_b:E_d\to E_{d+1}$ for its derivative operators and
$T_d(K):E_d\to E_d$ for the angular operators, real-linear in a skew
matrix $K$ and complex-linear in the field. Assume the derivative insertion
identities, the collinear angular Ward identities at every grade and arity,
and the commutation rule
\begin{equation}\label{eq:ward-analytic-commute}
 T_{d+1}(K)D_bv=D_bT_d(K)v+\sum_c K_{cb}D_cv.
\end{equation}
Signed coordinate symmetries, including reflection of the field labels,
are the ones already present in the hierarchy.

Fix a coordinate axis $a$. Its reflection-positive Hilbert space is
$\mathcal H_a$. For a field $v$ and a point $x$ with $x_a>0$, let
$\Psi_v(x)$ be the vector represented by the singleton insertion there.
Denote the reflected label by $v^{\dagger_a}$ and the coordinate reflection
by $\vartheta_a$. If
\[
 F_{v,w}(x)=S(v\text{ at }x,w\text{ at }0),
\]
the exact Gram identity is
\begin{equation}\label{eq:ward-analytic-gram}
 \langle\Psi_v(x),\Psi_w(y)\rangle
       =F_{w,v^{\dagger_a}}(y-\vartheta_a x).
\end{equation}
The inner product is linear in its second variable. Translation invariance
and the permutation symmetry of the hierarchy prove this identity directly
from the reflected-word definition of the Hilbert space.

\subsection{Smooth singleton vectors and the positive radial line}

For $b\ne a$, the curve $t\mapsto\Psi_v(x+te_b)$ is continuous by the
strongly continuous spatial translations on $\mathcal H_a$. Its derivative
is $\Psi_{D_bv}(x+te_b)$. To justify this as a Hilbert-space derivative,
first pair with any reflected-word vector. The scalar pairing is a smooth
separated-point moment, whose derivative is the stated derivative
insertion. The proposed derivative curve is itself continuous by spatial
translation. Its Bochner integral from $0$ to $t$ therefore has, against
every word vector, the same pairing as
$\Psi_v(x+te_b)-\Psi_v(x)$. Density of word vectors identifies these
vectors. The fundamental theorem of calculus now gives the strong
derivative. Iterating this argument with $D_bv,D_b^2v,\ldots$ proves
smoothness of the entire curve.

In the normal direction there is more: a singleton vector is analytic at
every positive normal coordinate. Choose an earlier strictly positive
normal coordinate. Nearby vectors are obtained from that earlier vector
by the heat semigroup at a strictly positive time. The holomorphic
positive-time extension of Section~\ref{sec:complex-heat} proves the
claim. Pairing with a fixed singleton and using
\eqref{eq:ward-analytic-gram} consequently shows that
\begin{equation}\label{eq:ward-radial-analytic}
 t\longmapsto F_{v,w}(te_a)
 \quad\hbox{is real analytic for }t>0.
\end{equation}
For example, near a specified $t_0>0$, place the fixed reflected singleton
at $t_0e_a/2$ and the variable singleton at $(t-t_0/2)e_a$.
Their separation is $te_a$ and both lie in the positive half-space.

\subsection{An analytic comparison with the same transverse jets}

Fix $b\ne a$ and the plane generator $K$ with
$Ke_a=e_b$, $Ke_b=-e_a$. For $u>0$ set
\[
 \rho(u,z)=\sqrt{u^2+z^2},\qquad
 \theta(u,z)=\arctan(z/u),
\]
\[
 Q_{v,w}(u,z)=
 F_{e^{-\theta(u,z)T(K)}v,\,e^{-\theta(u,z)T(K)}w}
                    (\rho(u,z)e_a).
\]
The exponentials act in the respective finite-dimensional field spaces.
Expand both labels in finite bases. Each summand is a product of analytic
matrix coefficients and a fixed radial moment from
\eqref{eq:ward-radial-analytic}, composed with the positive analytic
function $\rho$. Thus $Q_{v,w}$ is real analytic on $u>0$.

Define $\mathcal R=u\partial_z-z\partial_u$. Direct differentiation gives
$\mathcal R\theta=1$ and $\mathcal R\rho=0$. Consequently
\begin{equation}\label{eq:ward-comparator-pde}
 Q_{T(K)v,w}+Q_{v,T(K)w}+\mathcal RQ_{v,w}=0.
\end{equation}
Also $Q_{v,w}(u,0)=F_{v,w}(ue_a)$.
Let $G_{v,w}(u,z)=F_{v,w}(ue_a+ze_b)$ be the actual smooth two-point
function. Its corresponding expression
$G_{T(K)v,w}+G_{v,T(K)w}+\mathcal RG_{v,w}$ is the angular Ward expression
for these two insertions.

Every transverse derivative of that Ward expression vanishes at $(u,0)$.
Here is the reason this follows from the collinear identity for all labels.
Differentiate the angular expression in one insertion coordinate. The
new derivative of its linear spatial coefficient contributes
$\sum_c K_{cb}$ times the corresponding derivative insertion. Combining
this term with the differentiated spin insertion, using
\eqref{eq:ward-analytic-commute}, gives exactly the angular expression
with that label replaced by $D_bv$. All other terms are just differentiated
moment rows. Thus repeated coordinate differentiation only raises labels.
The assumed collinear identity applies to every resulting word. In
particular all $\partial_z^n$ of the two-point angular expression vanish
at every $(u,0)$ with $u>0$.

We show that $G$ and $Q$ have the same transverse jets, without assuming
that $G$ is analytic. Put $H=G-Q$, allowing every pair of labels in this
notation. Then $H_{v,w}(u,0)=0$ and all transverse jets of its angular
expression vanish. The elementary commutation formula is
\[
 \partial_z^n(\mathcal RH)
    =\mathcal R\partial_z^nH-n\partial_u\partial_z^{n-1}H
       \qquad(n\geq1).
\]
Assume inductively that all transverse jets through order $n$ of every
$H_{v,w}$ vanish along the positive $u$-axis. Apply $\partial_z^n$ to
its angular expression and set $z=0$. Both spin terms vanish by the
induction hypothesis. The last term in the displayed formula also
vanishes, because the preceding transverse jet is identically zero as a
function of $u$. We obtain
$u\partial_z^{n+1}H_{v,w}(u,0)=0$.
For $n=0$ the same conclusion follows directly without the last term.
Since $u>0$, induction proves
\begin{equation}\label{eq:ward-comparison-jets}
 \partial_z^nG_{v,w}(u,0)=\partial_z^nQ_{v,w}(u,0)
             \qquad(n\geq0).
\end{equation}
This equality of jets alone is not an analytic continuation statement.
The next Hilbert-space estimate supplies the required convergence.

\subsection{From scalar jets to a Hilbert analytic curve}

We use a useful elementary fact. Suppose $f:\mathbb R\to\mathcal H$ is
smooth and
\[
 \langle f(s),f(t)\rangle=A(t-s).
\]
Suppose the derivatives of $A$ at zero agree with those of a scalar
function analytic at zero. The power series of that function gives
constants $C,R>0$ with
$|A^{(n)}(0)|\leq C R^{-n}n!$ for every $n$. Differentiating the Gram
identity in both variables yields
\[
 \langle f^{(m)}(s),f^{(n)}(t)\rangle
                  =(-1)^m A^{(m+n)}(t-s).
\]
In particular, uniformly for all real $t$,
\[
 \|f^{(n)}(t)\|^2
    =|A^{(2n)}(0)|\leq C R^{-2n}(2n)!
    \leq C(2/R)^{2n}(n!)^2.
\]
The last inequality uses $\binom{2n}{n}\leq2^{2n}$.
Taylor's integral remainder in the Hilbert space is therefore at most
$\sqrt C(2|h|/R)^{N+1}$ after order $N$. It tends to zero for
$|h|<R/2$. The Taylor series converges in norm and equals $f$ there,
proving real analyticity at every point of the line.

Apply this fact to $f(t)=\Psi_v(x+te_b)$ for $x_a>0$ and $b\ne a$.
By \eqref{eq:ward-analytic-gram},
\[
 A(t-s)=F_{v,v^{\dagger_a}}
                  (2x_ae_a+(t-s)e_b).
\]
The smoothness was proved above. The analytic comparator is
$z\mapsto Q_{v,v^{\dagger_a}}(2x_a,z)$, and
\eqref{eq:ward-comparison-jets} gives precisely the required derivatives
at zero. Hence the singleton trajectory is analytic along every transverse
coordinate, as well as along the positive normal coordinate. Stationarity
of the Gram kernel is what makes the factorial bound uniform along the
whole transverse line.

\subsection{Moving an extreme insertion in an arbitrary word}

Consider a configuration of distinct points and one insertion whose
$a$-coordinate is strictly larger than those of all other insertions.
Choose a separating hyperplane $x_a=c$ strictly between them. Translate
it to zero and reflect the fixed spectator word into the positive
half-space. Reflection positivity then represents the moment, locally in
any coordinate motion of the extreme insertion, as
\[
 \langle\Phi,\Psi_v(x+te_b)\rangle
\]
with a fixed word vector $\Phi$. Strict separation persists in a
neighborhood of the chosen parameter. If $b=a$, positive-time analyticity
applies; if $b\ne a$, use the transverse result just proved. Applying the
continuous linear functional $\langle\Phi,\cdot\rangle$ proves
analyticity of the moment row. Permutation symmetry moves any selected
slot to the first slot, so no restriction on its position in the word is
needed. Spectators may have equal $a$-coordinates, provided their points
remain distinct.

For an insertion strictly smaller in coordinate $a$, use the hierarchy's
actual signed reflection in that axis, with its reflected labels. It
turns the insertion into a strict maximum. A motion in coordinate $a$
changes parameter $t$ to $-t$; a transverse motion leaves $t$ unchanged.
Both substitutions preserve analyticity. This reflection uses the already
constructed signed symmetry, not a rotation action yet to be proved.

Every angular Ward expression is a finite sum of such moment rows, with
labels either turned or raised once and coefficients affine in the moving
coordinate. It is therefore analytic wherever the moving point is a
strict extremum in some fixed coordinate. Along a closed segment with
this property at every point, these neighborhoods cover the segment.
This is the precise segment analyticity needed for the one-dimensional
identity theorem in the global Ward propagation argument.
